\documentclass[12pt,a4paper]{article}
\usepackage{ra_preamble}
\title{\textbf{The Causal Firewall: Substrate-Independent Physical Constraints\\
on the Origin and Distribution of Biological Complexity}\\[0.4em]
\large\normalfont Relational Actualism Paper~IV\\
\normalsize Submitted to \textit{International Journal of Astrobiology}}
\author{Joshua F.~Sandeman\\[0.3em]
\small Independent Researcher, Salem, Oregon\\
\small \texttt{sansuikyo@gmail.com} $\cdot$ \texttt{relationalactualism.org}}
\date{April 2026}

\begin{document}
\maketitle
\begin{abstract}
We derive substrate-independent physical constraints on the origin of biological complexity from Relational Actualism's causal graph dynamics, requiring no assumptions about chemistry, chirality, or aqueous solvents.  The Causal Firewall is the Erd\H{o}s--R\'enyi percolation transition at actualization density $\mu=1$ — the same threshold governing QCD confinement, galactic rotation curves, quantum error correction fault-tolerance, cosmological seed viability, and the RA actualization selectivity maximum.  Five independent physical scales are unified at a single dimensionless critical parameter.  The minimum viable complexity for self-replication is $N_{\min}=2$ causally connected actualization events — a non-Markovian structure with causal memory — identical to the minimum viable universe seed (Paper~III).  At biological temperatures ($T=300\,\mathrm{K}$), the universal actualization timescale formula (Paper~II) gives $\tau_d \approx 20\,\mathrm{fs}$, equal to the vibrational decoherence timescale; critically, this uses the RA on-shell selection criterion to identify the correct bath (optical phonons, not Debye acoustic modes) without additional assumptions.  We establish two substrate-independent biosignature criteria: F1 (redundant environmental recording of actualization outcomes) and F2 (non-relativistic actualization rate), deriving them directly from the BDG dynamics rather than from chemical or biological assumptions.  An explicit worked example traces the causal memory requirement from abstract graph topology through the BDG action to the minimum polymer length for a viable RNA self-replicator, bridging the formalism directly to prebiotic chemistry.
\end{abstract}

\tableofcontents

%───────────────────────────────────────────────────────────────
\section{Introduction: Dissolving the Contingency}

The origin of biological complexity is conventionally treated as a problem of prebiotic chemistry: which molecules assembled, in which solvent, at which temperature?  Different assumptions give different answers, and the entire inquiry is contingent on Earth-like conditions.

RA dissolves this contingency.  The question is not ``which chemistry?'' but ``what causal structure does self-replication require?''  The answer is $N \geq 2$ causally connected actualization events — a non-Markovian structure with causal memory.  This condition is derived from physics, not assumed from biology, and it applies regardless of the molecular substrate.

\textbf{Substrate independence is the central claim, and the text will return to it repeatedly.}  Every section of this paper addresses the same question from a different angle: whether a given condition can be stated without reference to carbon, water, chirality, or any chemical specificity.  Where we invoke specific chemistry (the RNA worked example, Section~\ref{sec:rna}), we do so only to confirm the abstract prediction, not to derive it.

%───────────────────────────────────────────────────────────────
\section{The Five-Scale $\mu=1$ Unification}

The single dimensionless parameter $\mu=1$ (Planck actualization density) governs five independent physical scales:

\begin{enumerate}
  \item \textbf{QCD confinement}: free quarks $\to$ color-confined hadrons at the Erd\H{o}s--R\'enyi percolation of the color-charged graph.
  \item \textbf{Galactic rotation curves}: transition from $d=4$ to $d=2$ effective graph geometry in sparse regions (Paper~III, Section~3.3).
  \item \textbf{Quantum error correction}: fault-tolerance threshold $p_{\mathrm{th}} = \alpha_{\mathrm{EM}} = 1/137$ (Paper~V, Section~4.1).
  \item \textbf{RA selectivity maximum}: D4U02 proves $\mu^* \approx 1.019$ \statusDR\ (Paper~I, Theorem~1).
  \item \textbf{Causal Firewall}: minimum actualization density for a self-replicating causal structure to maintain identity against quantum decoherence.
\end{enumerate}

These five appearances of $\mu=1$ are the same geometric fact — the Erd\H{o}s--R\'enyi percolation transition for the Poisson-CSG graph at Planck density — manifesting in different physical contexts.  They are not independent occurrences connected by analogy; they are the same theorem applied at different scales.

%───────────────────────────────────────────────────────────────
\section{The Correct Environmental Bath at $T=300\,\mathrm{K}$}\label{sec:bath}

\subsection{On-Shell Bath Identification}

The universal actualization timescale formula (Paper~II, Eq.~(3)) requires identifying which environmental modes are on-shell.  There are two candidate baths in biological systems:

\textbf{Optical phonon bath:} characteristic frequency $\omega_{\mathrm{opt}} \approx 3\times10^{13}\,\mathrm{rad/s}$, correlation time $\tau_c \approx 20\,\mathrm{fs}$, energy ratio $\hbar\omega_{\mathrm{opt}}/(k_BT) \approx 0.76$ at $T=300\,\mathrm{K}$.

\textbf{Debye acoustic bath:} characteristic frequency $\omega_D \approx 10^{11}\,\mathrm{rad/s}$, correlation time $\tau_c \approx 8.3\,\mathrm{ps}$, energy ratio $\hbar\omega_D/(k_BT) \approx 3\times10^{-3}$ at $T=300\,\mathrm{K}$.

The RA on-shell criterion selects a bath mode as active if $\hbar\omega/(k_BT) \geq \DeltaS \approx 0.6$.  Optical phonons satisfy this at $300\,\mathrm{K}$; Debye acoustic modes do not.  This selection is a \emph{derived} result from the threshold criterion, not an additional assumption.  Standard decoherence theory (Lindblad) treats both baths equally; RA differentiates them by the threshold.

\subsection{Derivation of $\tau_d$}

Applying the formula with the on-shell bath identified:
\begin{equation}
  \tau_d = \frac{\DeltaS}{\Gamma_{\mathrm{opt}}\cdot(\hbar\omega_{\mathrm{opt}}/k_BT)} = \frac{0.601}{3.93\times10^{13}\times 0.76} = 2.0\times10^{-14}\,\mathrm{s} = 20\,\mathrm{fs}
  \label{eq:taud}
\end{equation}
This equals $\tau_{\mathrm{vib}}$, the vibrational decoherence timescale.  The result has no free parameters: $\DeltaS$ is derived from the BDG model (Paper~I), $\Gamma_{\mathrm{opt}}$ is determined by temperature, and the bath identification follows from the on-shell criterion.

\subsection{Non-Markovian Correction}

For prebiotic chemical environments with reorganization energy $\lambda_{\mathrm{reorg}} \ll k_BT$ (slow-modulation limit), the Kubo line-shape theory gives a correction:
\begin{equation}
  \kappa = \frac{\sqrt{2\DeltaS}}{\mu_{\mathrm{NM}}}, \qquad N_{\min}^{\mathrm{NM}} = \left\lceil \frac{1}{\kappa}\right\rceil
  \label{eq:nonmarkov}
\end{equation}
For RNA-world conditions ($\lambda_{\mathrm{reorg}} \approx 0.1$--$0.3\,\mathrm{eV}$): $\kappa \approx 0.2$--$0.5$, giving $N_{\min}^{\mathrm{NM}} \approx 2$--$5$.  This is the range of minimum polymer lengths predicted by RA for the first self-replicating systems, consistent with known minimum RNA chain lengths for functional ribozymes ($\sim 10$--$30$ nucleotides for minimal hairpin ribozymes \cite{Joyce2018}).

%───────────────────────────────────────────────────────────────
\section{From Abstract Topology to Prebiotic Chemistry}\label{sec:rna}

\subsection{The Causal Memory Requirement}

The minimum viable complexity condition $N \geq 2$ connected actualization events is an \emph{abstract} statement about causal graph structure.  We now trace it to concrete molecular terms.

A self-replicating system requires:
\begin{enumerate}
  \item An actualization event $v_0$ encoding chemical information (a specific molecular bond forming, closing a ring, setting a stereocentre).
  \item A causally connected downstream event $v_1$ that uses the information from $v_0$ to reproduce the initial condition — i.e., a template-dependent reaction.
  \item A directed edge $v_0 \prec v_1$ (causal memory): the outcome of $v_0$ constrains the outcome of $v_1$.
\end{enumerate}
This is precisely $N=2$ with a directed edge, giving $V = \SBDG = 1 - 1 + 9 = 9$.  Without the directed edge ($N=2$ disconnected), $V=0$: the system cannot bootstrap.

\subsection{The RNA Worked Example}

Consider the minimal RNA self-replicating loop.  An RNA strand acts as a template for its complement; the complement then templates the original sequence.  In RA terms:
\begin{itemize}
  \item $v_0$: nucleotide addition to the growing strand — an on-shell phosphodiester bond formation ($\Delta S > \DeltaS$, irreversible in aqueous conditions).
  \item $v_1$: template-directed nucleotide addition to the complementary strand, causally constrained by the sequence established at $v_0$.
  \item Edge $v_0 \prec v_1$: the base-pairing specificity of the template is the causal connection.
\end{itemize}

The minimum strand length for this causal structure: a 2-nucleotide seed (the $N=2$ connected graph) is kinematically sufficient for self-complementarity but too short to fold stably.  The non-Markovian correction (Eq.~\ref{eq:nonmarkov}) with $\kappa \approx 0.3$ for aqueous RNA conditions gives $N_{\min}^{\mathrm{NM}} \approx 3$--$4$.  This corresponds to minimum RNA strand lengths of $\sim 6$--$8$ nucleotides — consistent with the known minimum for stable RNA hairpin formation ($\sim 4$--$8$ nucleotides).

\textbf{This is not a chemical derivation.}  We derived $N_{\min} = 2$ from the BDG action and corrected it for the non-Markovian bath environment.  The agreement with RNA-world minimum strand lengths is a \emph{confirmation} of the abstract prediction, not the basis for it.  A different chemical substrate (peptide nucleic acids, glycol nucleic acids, mineral-surface templates) with different $\kappa$ would give a different minimum length — but the same abstract condition $N\geq 2$ with directed edge.

\subsection{The Origin-of-Life Sandwich Bound}

Any proposed origin-of-life mechanism must produce a system whose RA assembly depth $A_{\mathrm{RA}}$ lies between:
\begin{itemize}
  \item Lower bound: $A_{\mathrm{RA}} \geq 2$ (non-Markovian causal memory with directed edge).
  \item Upper bound: $A_{\mathrm{RA}} \leq A_{\mathrm{max}}$ derived from the causal graph topology at the Planck scale.
\end{itemize}

The glycolysis pathway as a worked example: $A_{\mathrm{RA}} = 3$ (three non-trivial actualization depth levels, corresponding to the three main reaction stages — substrate phosphorylation, isomerization, and dephosphorylation).  Glycolysis satisfies $2 \leq 3 \leq A_{\mathrm{max}}$.

RNA-world, metabolism-first, and crystal-template scenarios can each be evaluated against this bound using biochemical databases (KEGG/BiGG).  The bound is substrate-independent: it constrains the causal depth structure, not the molecular identity.

%───────────────────────────────────────────────────────────────
\section{Substrate-Independent Biosignature Criteria}

\subsection{F1: Quantum Darwinism Redundancy}

\begin{definition}[Condition F1]
An environment satisfies F1 if it records the outcomes of actualization events redundantly enough to imprint them on macroscopic timescales — i.e., if the information content of actualization events is amplified into multiple environmental copies before the next actualization cycle.
\end{definition}

F1 makes no reference to carbon, water, molecular chirality, or any chemical substrate.  It is a condition on the information-carrying capacity of the environment.  Any medium satisfying the redundancy criterion qualifies: mineral surfaces, aerosol droplets, cosmic dust grains, ionic channels in exotic solvents.

\subsection{F2: Non-Relativistic Actualization}

\begin{definition}[Condition F2]
An environment satisfies F2 if the local effective temperature $T_{\mathrm{eff}}$ satisfies $T_{\mathrm{eff}} \ll T_U(a_{\mathrm{loc}})$, where $T_U = \hbar a/(2\pi c k_B)$ is the Unruh temperature at local acceleration $a_{\mathrm{loc}}$.
\end{definition}

In the extreme Rindler limit ($T_{\mathrm{eff}} \gg T_U$), $\Delta S_{\mathrm{per\;photon}} \to 0$ (by L05~\statusLV, Paper~II) and the actualization rate becomes uncontrolled.  A system swept past a gravitational horizon cannot support stable biological complexity.

\subsection{The Biosignature Prediction (P6)}

\begin{equation}
  \text{Life is found wherever F1 and F2 hold, regardless of chemistry.}
  \label{eq:P6}
\end{equation}

\textbf{Environments to search beyond the conventional habitable zone:}
\begin{itemize}
  \item Subsurface rock-hosted systems: F1 via mineral-water interfaces; F2 via suppressed acoustic modes at $T \sim 40$--$80\,^\circ\mathrm{C}$.
  \item High-pressure ices: F1 via ionic channels; F2 via low thermal noise.
  \item Interstellar grains: F1 marginal; F2 satisfied.
\end{itemize}

\textbf{Environments where life is prohibited:} extreme radiation environments (F1 destroyed), near stellar surfaces ($T_{\mathrm{eff}} \approx T_U$ near neutron star surfaces — F2 violated), inside active galactic nuclei.

%───────────────────────────────────────────────────────────────
\section{Discussion}

The Causal Firewall provides a precise, parameter-free account of what biological complexity requires physically.  F1 and F2 are physical theorems derived from the RA actualization criterion — they are not biochemical assumptions.  The five-scale $\mu=1$ unification demonstrates that the origin of life is not a special problem of chemistry but a particular manifestation of the same percolation physics governing QCD, gravitation, quantum computing, and the origin of universes.

The theory does not say which chemistry will produce life; it says \emph{what any chemistry must produce} — a directed causal connection between at least two actualization events.  The identification of which chemistry achieves this in a given environment is an application of the abstract constraint, not a derivation from it.

\section*{Acknowledgments}
The connection between the COSMO-SEED $N_{\min}=2$ result and the Causal Firewall was identified in collaboration with Claude (Anthropic).

\bibliographystyle{plain}
\begin{thebibliography}{99}
\bibitem{RA_P2} Sandeman, J.~F. (2026). Paper~II: Wave Function Collapse. Zenodo, doi:10.5281/zenodo.19362968.
\bibitem{RA_P3} Sandeman, J.~F. (2026). Paper~III: Cosmology. Zenodo, doi:10.5281/zenodo.19363017.
\bibitem{RA_P5} Sandeman, J.~F. (2026). Paper~V: The Complete Framework. Zenodo, doi:10.5281/zenodo.19363077.
\bibitem{Joyce2018} Joyce, G.~F., and Szostak, J.~W. (2018). Protocells and RNA self-replication. \textit{Cold Spring Harb.\ Perspect.\ Biol.}, 10, a034801.
\bibitem{Walker2013} Walker, S.~I., and Davies, P.~C.~W. (2013). The algorithmic origins of life. \textit{J.\ R.\ Soc.\ Interface}, 10, 20120869.
\bibitem{Cronin2016} Cronin, L., and Walker, S.~I. (2016). Beyond prebiotic chemistry. \textit{Science}, 352, 1174.
\bibitem{Zurek2009} Zurek, W.~H. (2009). Quantum Darwinism. \textit{Nature Physics}, 5, 181.
\bibitem{Baross1985} Baross, J.~A., and Hoffman, S.~E. (1985). Submarine hydrothermal vents. \textit{Origins of Life}, 15, 327.
\bibitem{Milgrom1983} Milgrom, M. (1983). A modification of the Newtonian dynamics. \textit{ApJ}, 270, 365.
\end{thebibliography}
\end{document}
